\documentclass[aspectratio=169,10pt,spanish]{beamer}

\usepackage[utf8]{inputenc}
\usepackage[T1]{fontenc}
\usepackage{lmodern}
\usepackage{helvet}
\renewcommand{\familydefault}{\sfdefault}
\usepackage[spanish,es-tabla,es-nodecimaldot]{babel}
\usepackage{tikz}
\usetikzlibrary{shapes.geometric,arrows.meta,positioning,calc,fit,backgrounds}
\usepackage{tcolorbox}
\tcbuselibrary{skins,breakable}
\usepackage{amsmath,amssymb}
\usepackage{colortbl}
\usepackage{booktabs}
\usepackage{tabularx}
\usepackage{array}
\usepackage{ragged2e}
\usepackage{hyperref}

\setbeamertemplate{navigation symbols}{}
\setbeamercolor{background canvas}{bg=white}
\setbeamersize{text margin left=0.55cm,text margin right=0.55cm}

% Paleta Blueprint utilizada en las presentaciones de Título y Objetivos.
\definecolor{bpNavy}{RGB}{28,43,54}
\definecolor{bpBlue}{RGB}{43,91,132}
\definecolor{bpGold}{RGB}{200,157,66}
\definecolor{bpRust}{RGB}{179,92,62}
\definecolor{bpDark}{RGB}{33,37,41}
\definecolor{bpLightGray}{RGB}{245,247,250}
\definecolor{bpGridLine}{RGB}{225,230,238}
\definecolor{bpGreen}{RGB}{46,125,50}
\definecolor{bpRed}{RGB}{198,40,40}

\setbeamertemplate{frametitle}{
    \begin{tikzpicture}[remember picture,overlay]
        \fill[bpNavy] (current page.north west) rectangle ([yshift=-1.05cm]current page.north east);
        \node[anchor=west,text=white,font=\bfseries\sffamily\normalsize]
            at ([xshift=0.6cm,yshift=-0.52cm]current page.north west) {\insertframetitle};
        \node[anchor=east,text=white!75!white,font=\tiny\sffamily]
            at ([xshift=-0.6cm,yshift=-0.52cm]current page.north east) {SEMINARIO DE INVESTIGACIÓN II};
    \end{tikzpicture}
    \vspace{0.25cm}
}

\setbeamertemplate{footline}{
    \begin{tikzpicture}[remember picture,overlay]
        \draw[bpGridLine,line width=0.8pt]
            ([yshift=0.5cm]current page.south west) -- ([yshift=0.5cm]current page.south east);
        \node[anchor=west,text=bpNavy!80!black,font=\tiny\sffamily]
            at ([xshift=0.6cm,yshift=0.25cm]current page.south west)
            {UNSAAC -- Ingeniería Informática y de Sistemas $\mid$ Br. Edmil Jampier Saire Bustamante};
        \node[anchor=east,text=bpNavy!80!black,font=\tiny\sffamily]
            at ([xshift=-0.6cm,yshift=0.25cm]current page.south east)
            {\insertframenumber{} / \inserttotalframenumber};
    \end{tikzpicture}
}

\newcommand{\fuente}[1]{\vspace{0.05cm}\hfill{\color{bpNavy!65}\tiny\textit{Seminario: diapositivas #1}}}
\newcommand{\ok}{\textcolor{bpGreen}{\bfseries\checkmark}}
\newcommand{\pendiente}{\textcolor{bpRust}{\bfseries\(\triangle\)}}

\begin{document}

% 1. Portada con el mismo tratamiento Blueprint de las presentaciones anteriores.
\begin{frame}[plain]
    \begin{tikzpicture}[remember picture,overlay]
        \draw[step=0.5cm,bpGridLine,line width=0.3pt]
            (current page.south west) grid (current page.north east);
        \draw[bpNavy,line width=1.4pt]
            ([xshift=0.5cm,yshift=0.5cm]current page.south west)
            rectangle ([xshift=-0.5cm,yshift=-0.5cm]current page.north east);
        \fill[bpNavy]
            ([xshift=0.5cm,yshift=-1.6cm]current page.north west)
            rectangle ([xshift=-0.5cm,yshift=-0.5cm]current page.north east);
        \node[anchor=center,text=white,font=\large\bfseries\sffamily,align=center,
            text width=12.8cm] at ([yshift=-1.02cm]current page.north)
            {INTRODUCCIÓN DEL PLAN DE TESIS};
    \end{tikzpicture}
    \par\vspace*{1.35cm}
    \centering
    {\color{bpNavy}\Large\bfseries Aplicación de los seis bloques}\\[0.2cm]
    {\color{bpDark}\normalsize Plan de tesis: arquitectura SLM--RAG con QLoRA para Quechua Collao}\\[0.45cm]
    \begin{tcolorbox}[width=12.4cm,colback=bpLightGray,colframe=bpNavy!65!black,
        arc=1.5mm,center]
        \centering\small\sffamily
        \textbf{Secuencia aplicada:}\\
        contexto \(\rightarrow\) tecnología y método \(\rightarrow\) vacío
        \(\rightarrow\) problema \(\rightarrow\) propuesta y justificación
        \(\rightarrow\) organización
    \end{tcolorbox}
    \vfill
    {\tiny\textbf{Tesista:} Br. Edmil Jampier Saire Bustamante
    \quad$\mid$\quad\textbf{Cátedra:} Msc. Raúl Huillca Huallparimachi\\
    UNSAAC -- Facultad de Ingeniería Eléctrica, Electrónica, Informática y Mecánica
    \quad$\mid$\quad Cusco, 2026}
\end{frame}

% 2. Mapa del proceso con trazabilidad a las 37 diapositivas de la fuente.
\begin{frame}{Ruta del seminario aplicada al plan de tesis}
    \begin{columns}[T,onlytextwidth]
        \begin{column}{0.32\textwidth}
            \begin{tcolorbox}[height=1.68cm,colback=bpRust!7,colframe=bpRust,arc=1.2mm,valign=center]
                \scriptsize\textbf{\color{bpRust}01 \quad Contexto}\\[1mm]
                \raggedright\tiny Campo, necesidad, dificultad y tema.\\[1mm]
                \tiny Fuente: 05--07
            \end{tcolorbox}
        \end{column}
        \begin{column}{0.32\textwidth}
            \begin{tcolorbox}[height=1.68cm,colback=bpGold!9,colframe=bpGold!90!black,arc=1.2mm,valign=center]
                \scriptsize\textbf{\color{bpGold!80!black}02 \quad Tecnología y método}\\[1mm]
                \raggedright\tiny Métodos, fundamento, evidencia y límites.\\[1mm]
                \tiny Fuente: 08--10
            \end{tcolorbox}
        \end{column}
        \begin{column}{0.32\textwidth}
            \begin{tcolorbox}[height=1.68cm,colback=bpBlue!7,colframe=bpBlue,arc=1.2mm,valign=center]
                \scriptsize\textbf{\color{bpBlue}03 \quad Vacío}\\[1mm]
                \raggedright\tiny Síntesis, contraste, limitación y delimitación.\\[1mm]
                \tiny Fuente: 11--16
            \end{tcolorbox}
        \end{column}
    \end{columns}
    \vspace{0.12cm}
    \begin{columns}[T,onlytextwidth]
        \begin{column}{0.32\textwidth}
            \begin{tcolorbox}[height=1.68cm,colback=bpRust!7,colframe=bpRust,arc=1.2mm,valign=center]
                \scriptsize\textbf{\color{bpRust}04 \quad Problema específico}\\[1mm]
                \raggedright\tiny Objeto, condiciones, comparación y criterios.\\[1mm]
                \tiny Fuente: 17--21
            \end{tcolorbox}
        \end{column}
        \begin{column}{0.32\textwidth}
            \begin{tcolorbox}[height=1.68cm,colback=bpBlue!7,colframe=bpBlue,arc=1.2mm,valign=center]
                \scriptsize\textbf{\color{bpBlue}05 \quad Propuesta y valor}\\[1mm]
                \raggedright\tiny Objetivo, justificación y evaluación.\\[1mm]
                \tiny Fuente: 21--32
            \end{tcolorbox}
        \end{column}
        \begin{column}{0.32\textwidth}
            \begin{tcolorbox}[height=1.68cm,colback=bpGreen!7,colframe=bpGreen,arc=1.2mm,valign=center]
                \scriptsize\textbf{\color{bpGreen}06 \quad Organización}\\[1mm]
                \raggedright\tiny Estructura real del documento de tesis.\\[1mm]
                \tiny Fuente: 33--34
            \end{tcolorbox}
        \end{column}
    \end{columns}
    \vspace{0.15cm}
    \begin{tcolorbox}[colback=bpNavy,colframe=bpNavy,arc=1mm,top=1mm,bottom=1mm]
        \centering\tiny\color{white}
        \textbf{Cobertura de la fuente:} 01--02 portada e índice; 03--04 arquitectura de seis bloques;
        05--34 desarrollo y ejemplos; 35--36 síntesis y lista de control; 37 encabezado de referencias
        sin referencias legibles. Las 37 imágenes y su extracción permanecen en la carpeta del proyecto.
    \end{tcolorbox}
    \fuente{03--04, 35--37}
\end{frame}

% 3. Bloque 1 aplicado.
\begin{frame}{Bloque 1: del campo general al tema específico}
    \begin{columns}[T,onlytextwidth]
        \begin{column}{0.49\textwidth}
            \begin{tcolorbox}[colback=bpLightGray,colframe=bpRust,title=\textbf{A. Campo}]
                \small El procesamiento del lenguaje natural permite analizar y generar textos mediante modelos computacionales.
            \end{tcolorbox}
            \begin{tcolorbox}[colback=white,colframe=bpGold!90!black,title=\textbf{B. Necesidad}]
                \small La utilidad de estas herramientas depende también de corpus y recursos lingüísticos que representen la variedad estudiada.
            \end{tcolorbox}
        \end{column}
        \begin{column}{0.49\textwidth}
            \begin{tcolorbox}[colback=white,colframe=bpRed,title=\textbf{C. Dificultad general}]
                \small Los recursos disponibles y las condiciones de ejecución pueden limitar la generación contextualizada en lenguas con menos recursos digitales.
            \end{tcolorbox}
            \begin{tcolorbox}[colback=bpBlue!6,colframe=bpBlue,title=\textbf{D. Delimitación}]
                \small El plan se centra en generar respuestas contextualizadas en Quechua Collao dentro de un entorno offline.
            \end{tcolorbox}
        \end{column}
    \end{columns}
    \vspace{0.05cm}
    \begin{tcolorbox}[colback=bpNavy,colframe=bpNavy,arc=1mm]
        \centering\scriptsize\color{white}
        \textbf{Transición redactada:} de la función del PLN y los recursos lingüísticos se pasa a la necesidad
        de evaluar generación contextualizada para Quechua Collao. Las afirmaciones sobre conectividad,
        cobertura digital o escasez deben apoyarse con fuentes del plan.
    \end{tcolorbox}
\end{frame}

% 4. Bloque 2: cuatro componentes para describir la tecnología y revisar evidencia.
\begin{frame}{Bloque 2: tecnología, fundamento y evidencia}
    \begin{columns}[T,onlytextwidth]
        \begin{column}{0.49\textwidth}
            \begin{tcolorbox}[colback=bpLightGray,colframe=bpBlue,title=\textbf{A--B. Métodos y función},fonttitle=\small\bfseries]
                \scriptsize
                \begin{itemize}
                    \item \textbf{SLM compacto:} produce la respuesta; el plan propone la familia Qwen2.5.
                    \item \textbf{RAG local:} recupera fragmentos de la base documental para aportar contexto.
                    \item \textbf{QLoRA:} adapta parámetros de manera eficiente; su aporte se debe medir.
                \end{itemize}
                \textbf{Aplicación:} integrar generación, recuperación local y adaptación para el dominio Quechua Collao.
            \end{tcolorbox}
        \end{column}
        \begin{column}{0.49\textwidth}
            \begin{tcolorbox}[colback=white,colframe=bpGold!90!black,title=\textbf{C--D. Evidencia y límites},fonttitle=\small\bfseries]
                \scriptsize
                \begin{itemize}
                    \item Orcotoma et al. (2025) y Duran y Silberztein (2026) aportan recursos lingüísticos.
                    \item Court y Elsner (2024) analizan contexto en traducción de quechua sureño al español.
                    \item Hevia et al. (2025) evalúan RAG offline en otro dominio y reportan dificultades con contexto irrelevante o extenso.
                \end{itemize}
                Esas tareas y condiciones no demuestran el comportamiento del sistema propuesto para Quechua Collao.
            \end{tcolorbox}
        \end{column}
    \end{columns}
    \begin{tcolorbox}[colback=bpNavy,colframe=bpNavy,arc=1mm,top=1.2mm,bottom=1.2mm]
        \centering\small\color{white}
        \textbf{Arquitectura a evaluar:}\quad
        SLM compacto \quad+\quad recuperación local (RAG) \quad+\quad adaptación (QLoRA)
    \end{tcolorbox}
    \fuente{08--10}
\end{frame}

% 5. Bloque 3: vacio delimitado a partir de antecedentes, no afirmado como hecho absoluto.
\begin{frame}{Bloque 3: construcción del vacío con evidencia}
    \vspace{-0.12cm}
    \begin{center}
    \scriptsize
    \begin{tabularx}{\textwidth}{>{\raggedright\arraybackslash}p{2.55cm} X X}
        \toprule
        \textbf{\color{bpNavy}Antecedente revisado} &
        \textbf{\color{bpNavy}Qué aporta} &
        \textbf{\color{bpNavy}Qué deja abierto para este plan}\\
        \midrule
        Corpus y recursos para Quechua Collao &
        Mejoran la disponibilidad de datos y recursos léxicos. &
        No evalúan por sí solos generación contextualizada mediante SLM--RAG--QLoRA.\\
        \addlinespace[1.5mm]
        Recuperación contextual en traducción &
        Examina el uso del contexto en una tarea de traducción de quechua sureño. &
        No equivale a generación de respuestas en Quechua Collao bajo inferencia offline.\\
        \addlinespace[1.5mm]
        RAG offline con modelos pequeños &
        Muestra retos de contexto en una aplicación de otro dominio. &
        Sus resultados no se trasladan automáticamente a datos ni tareas de esta tesis.\\
        \bottomrule
    \end{tabularx}
    \end{center}
    \vspace{0.08cm}
    \begin{tcolorbox}[colback=bpGold!12,colframe=bpGold!90!black,arc=1.2mm]
        \scriptsize
        \textbf{Vacío provisional:} falta determinar, con una revisión bibliográfica actualizada,
        cómo se comporta la combinación de modelo compacto, recuperación local y QLoRA al generar
        respuestas contextualizadas en Quechua Collao bajo restricciones offline.
        La novedad no se presenta como una afirmación universal.
    \end{tcolorbox}
    \fuente{11--16}
\end{frame}

% 6. Bloque 4: pregunta y anatomía.
\begin{frame}{Bloque 4: delimitación del problema específico}
    \begin{columns}[T,onlytextwidth]
        \begin{column}{0.24\textwidth}
            \begin{tcolorbox}[height=1.45cm,colback=bpBlue!7,colframe=bpBlue]
                \centering\scriptsize\textbf{\color{bpBlue}Objeto}\\
                Arquitectura SLM--RAG con QLoRA
            \end{tcolorbox}
        \end{column}
        \begin{column}{0.24\textwidth}
            \begin{tcolorbox}[height=1.45cm,colback=bpRust!7,colframe=bpRust]
                \centering\scriptsize\textbf{\color{bpRust}Contexto}\\
                Quechua Collao e inferencia offline
            \end{tcolorbox}
        \end{column}
        \begin{column}{0.24\textwidth}
            \begin{tcolorbox}[height=1.45cm,colback=bpGold!10,colframe=bpGold!90!black]
                \centering\scriptsize\textbf{\color{bpGold!80!black}Comparadores}\\
                SLM base y RAG estándar
            \end{tcolorbox}
        \end{column}
        \begin{column}{0.24\textwidth}
            \begin{tcolorbox}[height=1.45cm,colback=bpGreen!7,colframe=bpGreen]
                \centering\scriptsize\textbf{\color{bpGreen}Criterios}\\
                Fidelidad, anclaje y costo
            \end{tcolorbox}
        \end{column}
    \end{columns}
    \vspace{0.15cm}
    \begin{tcolorbox}[enhanced,colback=white,colframe=bpNavy,title=\textbf{Pregunta de trabajo propuesta},arc=1.5mm]
        \small\centering
        \textquestiondown De qué manera el diseño e implementación de una arquitectura SLM--RAG con QLoRA
        influye en la generación contextualizada en Quechua Collao frente al SLM base y al
        RAG estándar, bajo inferencia offline, considerando la fidelidad lingüística,
        el anclaje factual y el costo computacional?
    \end{tcolorbox}
    \begin{tcolorbox}[colback=bpLightGray,colframe=bpGridLine,arc=1mm,top=1mm,bottom=1mm]
        \centering\tiny
        La pregunta completa debe aparecer una sola vez en \textit{Formulación del problema}.
        En la introducción se anticipa el problema sin duplicarla.
    \end{tcolorbox}
    \fuente{17--21}
\end{frame}

% 7. Bloque 5: propuesta, objetivo y justificacion.
\begin{frame}{Bloque 5: propuesta, objetivo y justificación}
    \begin{tcolorbox}[colback=bpBlue!6,colframe=bpBlue,title=\textbf{Propuesta aplicada},fonttitle=\small\bfseries,top=0.8mm,bottom=0.8mm]
        \scriptsize Diseñar e implementar una arquitectura compacta que combine SLM,
        recuperación documental local y adaptación QLoRA; evaluar su comportamiento
        frente al modelo base y al RAG estándar.
    \end{tcolorbox}
    \begin{tcolorbox}[colback=bpLightGray,colframe=bpGold!90!black,title=\textbf{Objetivo general del plan, conservado},fonttitle=\small\bfseries,top=0.8mm,bottom=0.8mm]
        \scriptsize\centering
        Diseñar e implementar una arquitectura SLM--RAG con QLoRA para la generación
        contextualizada en Quechua Collao dentro de un entorno computacional offline.
    \end{tcolorbox}
    \begin{columns}[T,onlytextwidth]
        \begin{column}{0.32\textwidth}
            \begin{tcolorbox}[colback=white,colframe=bpBlue,title=\textbf{Ciencia},fonttitle=\scriptsize\bfseries,top=0.8mm,bottom=0.8mm]
                \tiny Evidencia sobre recuperación y adaptación para Quechua Collao.
            \end{tcolorbox}
        \end{column}
        \begin{column}{0.32\textwidth}
            \begin{tcolorbox}[colback=white,colframe=bpGold!90!black,title=\textbf{Método},fonttitle=\scriptsize\bfseries,top=0.8mm,bottom=0.8mm]
                \tiny Comparación reproducible entre configuraciones bajo condiciones definidas.
            \end{tcolorbox}
        \end{column}
        \begin{column}{0.32\textwidth}
            \begin{tcolorbox}[colback=white,colframe=bpRust,title=\textbf{Tecnología},fonttitle=\scriptsize\bfseries,top=0.8mm,bottom=0.8mm]
                \tiny Capacidades y límites de ejecución local en el entorno planteado.
            \end{tcolorbox}
        \end{column}
    \end{columns}
    \begin{tcolorbox}[colback=bpNavy,colframe=bpNavy,arc=1mm,top=0.8mm,bottom=0.8mm]
        \centering\tiny\color{white}
        En la introducción se anticipan comparación y criterios; parámetros y procedimientos quedan en el capítulo metodológico.
    \end{tcolorbox}
\end{frame}

% 8. Bloque 6: usar la estructura real del plan, no el ejemplo de cinco capitulos.
\begin{frame}{Bloque 6: organización real de la tesis}
    \begin{columns}[T,onlytextwidth]
        \begin{column}{0.32\textwidth}
            \begin{tcolorbox}[height=1.65cm,colback=bpBlue!7,colframe=bpBlue]
                \centering\small\textbf{\color{bpBlue}I. Introducción}\\[1mm]
                \scriptsize Problema, objetivos, justificación y alcance.
            \end{tcolorbox}
        \end{column}
        \begin{column}{0.32\textwidth}
            \begin{tcolorbox}[height=1.65cm,colback=bpGold!9,colframe=bpGold!90!black]
                \centering\small\textbf{\color{bpGold!80!black}II. Marco teórico}\\[1mm]
                \scriptsize Fundamentos y estado del arte.
            \end{tcolorbox}
        \end{column}
        \begin{column}{0.32\textwidth}
            \begin{tcolorbox}[height=1.65cm,colback=bpRust!7,colframe=bpRust]
                \centering\small\textbf{\color{bpRust}III. Metodología}\\[1mm]
                \scriptsize Diseño de la arquitectura y evaluación.
            \end{tcolorbox}
        \end{column}
    \end{columns}
    \vspace{0.16cm}
    \begin{columns}[T,onlytextwidth]
        \begin{column}{0.32\textwidth}
            \begin{tcolorbox}[height=1.65cm,colback=bpBlue!7,colframe=bpBlue]
                \centering\small\textbf{\color{bpBlue}IV. Implementación}\\[1mm]
                \scriptsize Desarrollo computacional del sistema.
            \end{tcolorbox}
        \end{column}
        \begin{column}{0.32\textwidth}
            \begin{tcolorbox}[height=1.65cm,colback=bpGold!9,colframe=bpGold!90!black]
                \centering\small\textbf{\color{bpGold!80!black}V. Resultados}\\[1mm]
                \scriptsize Resultados y evaluación experimental.
            \end{tcolorbox}
        \end{column}
        \begin{column}{0.32\textwidth}
            \begin{tcolorbox}[height=1.65cm,colback=bpGreen!7,colframe=bpGreen]
                \centering\small\textbf{\color{bpGreen}VI. Cierre}\\[1mm]
                \scriptsize Conclusiones y trabajos futuros.
            \end{tcolorbox}
        \end{column}
    \end{columns}
    \vspace{0.16cm}
    \begin{tcolorbox}[colback=bpNavy,colframe=bpNavy,arc=1mm,top=1mm,bottom=1mm]
        \centering\tiny\color{white}
        La secuencia se tomó de \texttt{partes/06cierre.tex}. El seminario ilustra cinco capítulos;
        para esta aplicación se conserva la estructura de seis capítulos del plan.
    \end{tcolorbox}
    \fuente{33--34}
\end{frame}

% 9. Síntesis narrativa: la cadena argumental aplicada en la introducción.
\begin{frame}{La introducción integrada: la cadena argumental aplicada}
    \begin{columns}[T,onlytextwidth]
        \begin{column}{0.49\textwidth}
            \begin{tcolorbox}[colback=bpRust!5,colframe=bpRust,title=\textbf{1. Contexto},fonttitle=\small\bfseries,top=0.8mm,bottom=0.8mm]
                \scriptsize El PLN procesa y genera texto; su aplicación depende de recursos que representen Quechua Collao.
            \end{tcolorbox}
            \begin{tcolorbox}[colback=bpGold!7,colframe=bpGold!90!black,title=\textbf{2. Tecnología y evidencia},fonttitle=\small\bfseries,top=0.8mm,bottom=0.8mm]
                \scriptsize El plan combina un modelo compacto, RAG local y QLoRA; los antecedentes estudian contextos diferentes.
            \end{tcolorbox}
            \begin{tcolorbox}[colback=bpBlue!5,colframe=bpBlue,title=\textbf{3. Vacío provisional},fonttitle=\small\bfseries,top=0.8mm,bottom=0.8mm]
                \scriptsize Se evaluará la combinación para generación en Collao y en entorno offline; la revisión bibliográfica debe confirmar el vacío.
            \end{tcolorbox}
        \end{column}
        \begin{column}{0.49\textwidth}
            \begin{tcolorbox}[colback=bpRust!5,colframe=bpRust,title=\textbf{4. Problema},fonttitle=\small\bfseries,top=0.8mm,bottom=0.8mm]
                \scriptsize Se comparan SLM base, RAG estándar y configuración integrada según fidelidad, anclaje y costo.
            \end{tcolorbox}
            \begin{tcolorbox}[colback=bpBlue!5,colframe=bpBlue,title=\textbf{5. Propuesta y valor},fonttitle=\small\bfseries,top=0.8mm,bottom=0.8mm]
                \scriptsize Se diseñará e implementará la arquitectura para generar texto contextualizado en Quechua Collao offline.
            \end{tcolorbox}
            \begin{tcolorbox}[colback=bpGreen!5,colframe=bpGreen,title=\textbf{6. Organización},fonttitle=\small\bfseries,top=0.8mm,bottom=0.8mm]
                \scriptsize La tesis sigue seis capítulos: introducción, teoría, metodología, implementación, resultados y conclusiones.
            \end{tcolorbox}
        \end{column}
    \end{columns}
\end{frame}

% 10. Panel de revisión para la integración posterior.
\begin{frame}{Panel de revisión antes de integrar el borrador}
    \begin{columns}[T,onlytextwidth]
        \begin{column}{0.49\textwidth}
            \begin{tcolorbox}[enhanced,colback=bpGreen!5,colframe=bpGreen,
                title=\textbf{Aplicación realizada},arc=1.5mm]
                \scriptsize
                \begin{tabular}{@{}p{0.35cm}p{5.55cm}@{}}
                    \ok & Se recorren los seis bloques del seminario y se aplican al tema del plan.\\[1mm]
                    \ok & Se asigna trazabilidad a las diapositivas fuente 01--37.\\[1mm]
                    \ok & Se usa la estructura real de seis capítulos del plan.\\[1mm]
                    \ok & Los ejemplos QFT/CKKS se mantienen como ejemplos de clase, no como tema de tesis.\\[1mm]
                    \ok & El plan original permanece intacto; el fragmento se entrega aparte para revisión.
                \end{tabular}
            \end{tcolorbox}
        \end{column}
        \begin{column}{0.49\textwidth}
            \begin{tcolorbox}[enhanced,colback=bpGold!8,colframe=bpRust,
                title=\textbf{Pendiente de validar en el plan},arc=1.5mm]
                \scriptsize
                \begin{tabular}{@{}p{0.35cm}p{5.55cm}@{}}
                    \pendiente & Confirmar el vacío con una búsqueda bibliográfica actualizada.\\[1mm]
                    \pendiente & Corregir o retirar \texttt{Chen2024}: su URL y su título local no corresponden.\\[1mm]
                    \pendiente & Sustentar con fuente la cifra de fertilidad de tokens cercana a \(3.4\times\), o quitarla.\\[1mm]
                    \pendiente & Revisar coherencia entre pregunta, objetivo general, objetivos específicos y métricas.
                \end{tabular}
            \end{tcolorbox}
        \end{column}
    \end{columns}
    \vspace{0.12cm}
    \begin{tcolorbox}[colback=bpNavy,colframe=bpNavy,arc=1mm,top=1mm,bottom=1mm]
        \centering\tiny\color{white}
        El borrador es una propuesta de redacción, no una afirmación de resultados ni una modificación automática del plan.
    \end{tcolorbox}
\end{frame}

% 11. Referencias citadas en la aplicación.
\begin{frame}{Fuentes citadas en la aplicación}
    \begin{columns}[T,onlytextwidth]
        \begin{column}{0.49\textwidth}
            \begin{tcolorbox}[colback=bpLightGray,colframe=bpBlue,arc=1.2mm]
                \scriptsize
                \textbf{Orcotoma Mormontoy et al. (2025).}
                \textit{Towards the Creation of a Collao Quechua--Spanish Parallel Corpus Using Optical Character Recognition.}
                \href{https://aclanthology.org/2025.globalnlp-1.1/}{ACL Anthology}
            \end{tcolorbox}
            \begin{tcolorbox}[colback=white,colframe=bpGold!90!black,arc=1.2mm]
                \scriptsize
                \textbf{Duran y Silberztein (2026).}
                \textit{Semi-Automatic Construction of a Quechua--Spanish Dictionary.}
                \href{https://aclanthology.org/2026.loresmt-1.9/}{ACL Anthology}
            \end{tcolorbox}
            \begin{tcolorbox}[colback=bpLightGray,colframe=bpRust,arc=1.2mm]
                \scriptsize
                \textbf{Court y Elsner (2024).}
                \textit{Shortcomings of LLMs for Low-Resource Translation: Retrieval and Understanding Are Both the Problem.}
                \href{https://aclanthology.org/2024.wmt-1.125/}{ACL Anthology}
            \end{tcolorbox}
            \begin{tcolorbox}[colback=white,colframe=bpBlue,arc=1.2mm]
                \scriptsize
                \textbf{Hevia et al. (2025).}
                \textit{Towards an Efficient, Customizable, and Accessible AI Tutor.}
                \href{https://arxiv.org/abs/2510.06255}{arXiv:2510.06255}
            \end{tcolorbox}
        \end{column}
        \begin{column}{0.49\textwidth}
            \begin{tcolorbox}[colback=bpLightGray,colframe=bpGold!90!black,arc=1.2mm]
                \scriptsize
                \textbf{Yang et al. (2024).}
                \textit{Qwen2.5 Technical Report.}
                \href{https://arxiv.org/abs/2412.15115}{arXiv:2412.15115}
            \end{tcolorbox}
            \begin{tcolorbox}[colback=white,colframe=bpBlue,arc=1.2mm]
                \scriptsize
                \textbf{Dettmers et al. (2023).}
                \textit{QLoRA: Efficient Finetuning of Quantized Language Models.}
                \href{https://arxiv.org/abs/2305.14314}{arXiv:2305.14314}
            \end{tcolorbox}
            \begin{tcolorbox}[colback=bpLightGray,colframe=bpRust,arc=1.2mm]
                \scriptsize
                \textbf{Xu et al. (2024).}
                \textit{On-Device Language Models: A Comprehensive Review.}
                \href{https://arxiv.org/abs/2409.00088}{arXiv:2409.00088}
            \end{tcolorbox}
            \begin{tcolorbox}[colback=bpRed!5,colframe=bpRed,arc=1.2mm]
                \tiny
                La diapositiva 37 del seminario solo muestra el encabezado de referencias.
                Estas fuentes proceden del borrador de aplicación y requieren revisión bibliográfica
                antes de incorporarse al documento final.
            \end{tcolorbox}
        \end{column}
    \end{columns}
    \fuente{37 y referencias del borrador}
\end{frame}

\end{document}
